% Faithful editable transcription of official CCC2024 article19, model pp.19:12--19:14 and protocol pp.19:15--19:16.
\section*{Author model and preliminaries}
We use the following standard result about the existence of error correcting codes.
\paragraph{Lemma 4.} Let $\delta>0$ and define $K_0=\lceil10/\delta^2\rceil$. For all $n>0$, there exists a function $\operatorname{ECC}_{n,\delta}:\{0,1\}^n\to\{0,1\}^{K_0 n}$ such that for all $s\ne t\in\{0,1\}^n$, we have
\[\Delta(\operatorname{ECC}_{n,\delta}(s),\operatorname{ECC}_{n,\delta}(t))>\left(\frac12-\delta\right)\cdot K_0 n.\]
Our communication model is the multi-party adversarial broadcast channel. Throughout this paper, we use $n$ to denote the number of parties. An $n$-party protocol in this model is defined by a tuple:
\[\Pi=\left(\{X^{(i)}\}_{i\in[n]},Y,T,\sigma,\{M_j\}_{j\in[T]},\operatorname{out}\right),\]
where
\begin{enumerate}
\item for all $i\in[n]$, $X^{(i)}$ is set of inputs of party $i$. We also define $X=X^{(1)}\times\cdots\times X^{(n)}$.
\item $Y$ is set of outputs of a protocol.
\item $T=\|\Pi\|\in\mathbb N$ is the length (number of rounds) in the protocol.
\item $\sigma\in[n]^T$ is a vector indicating for all rounds $j\in[T]$, which is the (unique) party scheduled to speak in round $j$.
\item for all $j\in[T]$, $M_j:X^{\sigma_j}\times\{0,1\}^{j-1}\times(\{0,1\}^*)^j\to\{0,1\}$ is the message function used by party $\sigma_j$ that uses his input, the bits he received in the first $j-1$ rounds, and the randomness sampled in the first $j$ rounds to output the bit he will broadcast in round $j$.
\item $\operatorname{out}:\{0,1\}^T\to Y$ is the function that the parties use to compute the output of the protocol based on the bits they receive.
\end{enumerate}
We will omit $X^{(i)}$ and $Y$ when they are clear from context.
\paragraph{Adversaries.} Let $\Pi$ be a protocol as above. An adversary for $\Pi$ is defined by a tuple $\operatorname{Adv}=(\operatorname{Adv}_{i,j})_{i\in[n],j\in[T]}$, where $\operatorname{Adv}_{i,j}:X\times(\{0,1\}^*)^j\to\{0,1\}$. Here, for all $i\in[n]$ and $j\in[T]$, the function $\operatorname{Adv}_{i,j}$ takes as input the inputs of all the parties and the randomness sampled in the first $j$ rounds and outputs 1 if he wants to flip (corrupt) the bit party $i$ receives in round $j$, and outputs 0 otherwise. Our formulation thus, captures adversaries that have knowledge of all the parties' inputs and the randomness they sampled so far, but are unaware of the randomness they will sample in the future.
\paragraph{Protocol execution.} We next describe the execution of a protocol $\Pi$ in the presence of adversary $\operatorname{Adv}$: Each party $i\in[n]$ starts with an input $x^{(i)}\in X^{(i)}$. Let $x=(x^{(1)},\ldots,x^{(n)})$. The execution takes place in $T$ rounds, maintaining the invariant that for all parties $i\in[n]$ and all rounds $j\in[T]$, party $i$ has a transcript $\pi_{<j}^{(i)}$ before the execution of round $j$. In each round $j\in[T]$, the parties first sample a shared random string $r_j\in\{0,1\}^*$. The player $\sigma_j$ then computes $\pi_j=M_j(x^{(\sigma_j)},\pi_{<j}^{(\sigma_j)},r_{\le j})$ and broadcasts it over the channel. All parties $i\in[n]$ then receive a (potentially corrupted) bit $\pi_j^{(i)}=\pi_j\oplus\operatorname{Adv}_{i,j}(x,r_{\le j})$ and append it to $\pi_{<j}^{(i)}$ to get $\pi_{\le j}^{(i)}$. At the end of the protocol, all parties $i\in[n]$ output $\operatorname{out}(\pi_{\le T}^{(i)})$. We say that an execution is noiseless if $\operatorname{Adv}_{i,j}$ always outputs 0, and we call this adversary the noiseless adversary.
Note that when $\Pi$ is executed in the presence of $\operatorname{Adv}$, the output of any party $i\in[n]$ is determined by the parties' inputs $x=(x^{(1)},\ldots,x^{(n)})$ and the sampled randomness $r_{\le T}$. Owing to this, we denote it using the notation $\Pi_{\operatorname{Adv},i}(x,r_{\le T})$. We will omit writing $\operatorname{Adv},i$ when the adversary is noiseless as in this case, all players compute the same transcripts, and therefore, also the same output. We also omit $r_{\le T}$ from our notation when talking about deterministic protocols. Also, as mentioned above, we define our adversaries to have complete knowledge of the inputs of the parties and the randomness they sampled so far but they are not aware of any future randomness the parties might have. Due to their knowledge of the randomness sampled so far (and the inputs), the adversaries can also compute all the bits sent and received over the channel so far, and we do not explicitly include this in our notation. Moreover, as our upper bound is deterministic, we will only need this assumption in our lower bound, and it only makes our result stronger.
Next, note that, as defined, the output function $\operatorname{out}$ is the same for all parties and depends only on the transcript of the protocol and not on the inputs of the parties. This is without loss of generality, as we can always extend the protocol so that one of the parties computes and sends the output over the channel (using an error correcting code) and all the parties then decode it from the transcript. Finally, note that our definition allows us to easily measure the number of bits corrupted by the adversary. As we are interested only in adversaries that do not corrupt too many bits sent or received by any given party, we define, for all $i\in[n]$, the set $\sigma^{-1}(i)=\{j\in[T]:\sigma(j)=i\}$ to contain rounds where party $i$ broadcast and also define:
\paragraph{Definition 5.} Let $\Pi$ be a protocol and $\operatorname{Adv}$ be an adversary for $\Pi$. Let $\theta>0$. We say that $\operatorname{Adv}$ is $\theta$-limited, if for all $x=(x^{(1)},\ldots,x^{(n)})$, all $r_{\le T}$, and all $i\in[n]$, we have:
\begin{align*}\sum_{j\in[T]}\operatorname{Adv}_{i,j}(x,r_{\le j})&\le\theta T,\\\sum_{j\in\sigma^{-1}(i)}\sum_{i'\in[n]}\operatorname{Adv}_{i',j}(x,r_{\le j})&\le\theta n\cdot|\sigma^{-1}(i)|.\end{align*}
Let $\Pi$ be a protocol as above and $f:X\to Y$ be a (possibly partial) function. Let $\theta,p>0$. We say that the protocol $\Pi$ computes $f$ with probability $p$ resilient to $\theta$-adversarial noise if for all inputs $x=(x^{(1)},\ldots,x^{(n)})$ in the domain of $f$ and all $\theta$-limited adversaries, we have:
\[\Pr(\forall i\in[n]:\Pi_{\operatorname{Adv},i}(x,r_{\le T})=f(x))\ge p.\]
We omit writing ``resilient to'' when $\theta=0$. As there is only one adversary that is 0-limited, we will also omit $\operatorname{Adv}$ from the subscript in this case.
We consider the $n$-party Information Dissemination function, denoted $\operatorname{ID}_n$, that simply outputs its $n$-sized tuple of arguments. That is, $\operatorname{ID}_n:\{0,1\}^n\to\{0,1\}^n$, where $\operatorname{ID}_n(x_1,\ldots,x_n)=(x_1,\ldots,x_n)$.
% STATEMENT_BEGIN
\paragraph{Theorem 8.} For all $\theta<1/40$, the protocol in Algorithm 1 solves $\operatorname{ID}_n$ resilient to $\theta$-adversarial noise.
% STATEMENT_END
% PROOF_BEGIN
\section*{4 Our Information Dissemination Protocol}
In order to demonstrate that our result in Theorem 6 is tight, we provide a simple algorithm with length $O(n^2)$, which is resilient to $\theta$-adversarial noise, for all $\theta<1/40$.
\paragraph{High level description.} To summarize, the idea is to proceed in two phases. In the first phase, every player says their input bit $O(n)$ times. Each other player then takes the majority value of what they heard. Thus, each player now has a guess for each player's input. Then, in the second phase, every player encodes the string of guesses they have using an error-correcting code, and broadcasts the result. Each player then takes all the error-correcitng codes they've received, decodes them, and sets their guess for each player's input to be the majority value among the guesses they've decoded.
The reason this works is that in order to corrupt a player's output, the adversary needs to either corrupt a lot of the error-correcting codes received by that player, or they need to corrupt a lot of players' guesses about some specific party's input. In either case, the adversary ends up corrupting too many rounds of communication, thus showing that a $\theta$-limited adversary cannot possibly corrupt even a single player's output, exactly as desired.
\paragraph{The formal protocol.} Our protocol is given in Algorithm 1. It uses an error correcting code $\operatorname{ECC}_{n,\delta}$ as promised by Lemma 4, where $\delta=1/10$. We use $K_0$ as given in that lemma.
\paragraph{Algorithm 1. The algorithm computing $\operatorname{ID}_n$.}
Input: Each party $k\in[n]$ has an input $x_k\in\{0,1\}$.
Output: Each party $i\in[n]$ outputs a $\hat x_{i,1},\ldots,\hat x_{i,n}$, such that $\hat x_{i,k}=x_k$ for all $k\in[n]$.
\begin{enumerate}
\item For $k\in[n]$ do
\item \quad Party $k$ broadcasts $x_k$ $K_0 n$ times.
\item \quad Each party $j\in[n]$ sets $y_{j,k}$ equal to the majority value they received in the previous $K_0 n$ broadcasts.
\item End for
\item For $j\in[n]$ do
\item \quad Party $j$ computes and broadcasts $\operatorname{ECC}_{n,\delta}(y_{j,1},\ldots,y_{j,n})$. This takes $K_0 n$ broadcasts.
\item \quad Each party $i\in[n]$ sets $\hat y_{i,j,1},\ldots,\hat y_{i,j,n}$ to the minimize the distance of $\operatorname{ECC}_{n,\delta}(\hat y_{i,j,1},\ldots,\hat y_{i,j,n})$ to the messages they received in the previous $K_0 n$ broadcasts.
\item End for
\item Each party $i\in[n]$ sets $\hat x_{i,k}$ to be the majority value between $\hat y_{i,1,k},\ldots,\hat y_{i,n,k}$ for all $k\in[n]$.
\item Each party $i\in[n]$ outputs $\hat x_{i,1},\ldots,\hat x_{i,n}$.
\end{enumerate}
\begin{proof}[Proof of Theorem 8]
We begin by noting that this protocol is deterministic. As such, the behavior of the protocol is completely determined by the inputs to the players and the adversary $\operatorname{Adv}$, and that, for simplicity, adversaries can be assumed to just be a function of the inputs $x_1,\ldots,x_n$.
Fix some adversary $\operatorname{Adv}$. We wish to demonstrate that if there exists some input $x_1,\ldots,x_n$ such that executing the protocol in Algorithm 1 against $\operatorname{Adv}$ on inputs $x_1,\ldots,x_n$ results in a player outputting an incorrect output, then the adversary $\operatorname{Adv}$ is not $\theta$-limited.
To see this, suppose that a player outputs some incorrect value. In particular, suppose that $\hat x_{i,k}\ne x_k$ for some $i\in[n]$ and $k\in[n]$. That means that the majority value among $\hat y_{i,1,k},\ldots,\hat y_{i,n,k}$ was not $x_k$.
That implies that there must exist some set $S\subseteq[n]$ such that $|S|\ge n/4$, and that one of the two following conditions must hold true:
\begin{enumerate}\item For all $j\in S$, $\hat y_{i,j,k}\ne y_{j,k}$, or\item for all $j\in S$, $y_{j,k}\ne x_k$.\end{enumerate}
We claim that in order for either of these cases to occur, the adversary $\operatorname{Adv}$ must not be $\theta$-limited. First suppose that for all $j\in S$, $\hat y_{i,j,k}\ne y_{j,k}$. That implies that for each $j\in S$, during iteration $j$ of the loop at Line 5, at least $0.2K_0 n$ of the bits sent by player $j$ are received incorrectly by player $i$, by the properties promised about $\operatorname{ECC}_{n,\delta}$ in Lemma 4. That implies that over all the iterations of the loop at Line 5, player $i$ hears at least $\frac{0.2}4 K_0 n^2=\frac1{20}K_0 n^2$ messages incorrectly. That means that at least $1/40$ of all $2K_0 n^2$ messages sent during the protocol are misheard by player $i$. This, thus, shows that $\operatorname{Adv}$ is not $\theta$-limited.
On the other hand, suppose that for all $j\in S$, $y_{j,k}\ne x_k$. That implies that for all $j\in S$, during iteration $k$ of the loop at Line 1, at least $0.5K_0 n$ of the bits sent by player $k$ are received incorrectly by player $j$. That implies that there are at least $\frac{0.5}4 K_0 n^2=1/8K_0 n^2$ corruptions in messages from player $k$ to other players, of a total of $2K_0 n^2$ messages received in rounds during which this player broadcasts. This, thus, shows that $\operatorname{Adv}$ is not $\theta$-limited.
\end{proof}
% PROOF_END
